\documentclass[12pt,a4paper]{article}
\usepackage[utf8]{inputenc}
\usepackage{amsmath, amssymb, bm}
\usepackage{algorithm}
\usepackage{algorithmic}
\usepackage{graphicx}
\usepackage{geometry}
\usepackage{caption}
\usepackage{cite}
\geometry{margin=1in}

% 启用中文支持
\usepackage{ctex}

% 将 hyperref 放置在最后，并显式配置内部链接和颜色
\usepackage[colorlinks=true, linkcolor=blue, anchorcolor=blue, citecolor=blue, urlcolor=blue, bookmarksnumbered=true]{hyperref}

\title{\textbf{基于“人类在环”贝叶斯优化的光催化 \\ CO$_2$ 还原多产物预测与材料设计架构}}
\author{Jiayu Wu}
\date{\today}

\begin{document}
\maketitle

\begin{abstract}
光催化二氧化碳还原反应 (CO$_2$RR) 的材料探索受限于多物理化学属性相互交织所带来的高维搜索空间，叠加实验试错的巨额时间成本，传统经验式寻优极为低效。本文提出并实现了针对该问题量身定制的“人类在环”贝叶斯优化（Human-in-the-Loop Bayesian Optimization, HITL-BO）计算框架。该框架支持包含 CO、HCOOH、CH$_4$ 在内的多种 CO$_2$ 还原反应产物的收率最优化。针对多尺度物理问题带来的“维度灾难”，体系创造性地将随机森林（Random Forest）降维前置，与具有自动关联性确定（ARD）机制的 Matérn 高斯过程核代理模型相融合，在保持贝叶斯推断严格性的前提下克服了特征冗余造成的模型光滑化。我们还引入了带噪声期望提升（qNEI）和人类干预策略（Manual Override），使得实验专家能够在自动推荐引擎和化学直觉间取得平衡。研究成果涵盖了由特征拆分、代理边界标定到数据闭环回填的详述及严谨的算法论断，充分借鉴了本领域的经典算法基石。

\textbf{关键词}：光催化；CO$_2$ 还原；贝叶斯优化；高斯过程；自动关联性决定；人类在环设计
\end{abstract}

\vspace{2em}
\tableofcontents
\newpage

\section{引言 (Introduction)}

光催化二氧化碳还原成为应对气候变化、将 CO$_2$ 转化为高价值化学品（如 CO、甲酸、甲烷甚至多碳醇类）的革新性道路之一\cite{nature2025, jacs2024}。然而，在实际研究与工业运用中，催化剂性能（收率法拉第效率）呈高度复杂的非凸变化。这些特性受到多达数十个参数的共同控制（如能带间隙、描述符 A\_pI、反应位点等），在如此高维度和实验测量存在固有噪音的影响下，催化系统表现出显著的非线性行为。以往依赖高通量筛选 (HTS)\cite{nature2025} 与全穷举实验的方法容易受到“维数灾难 (Curse of Dimensionality)”的桎梏。

应对昂贵评估函数的典型高效最优化范式当属贝叶斯优化 (Bayesian Optimization, BO) \cite{bo2012}。尽管标准 BO 理论对黑盒导数的评估是严密自洽的，但在真实的化学实验室设定里有明显鸿沟。首先是量纲差异下多元映射的不稳定性\cite{turbo2019}，其次则是不可执行点（物理配方不相容）引发的失效。因此，本报告构建的主体系统对 BoTorch 及 GPyTorch 等库深度魔改优化，将理论 GP-UCB (Gaussian Process - Upper Confidence Bound)\cite{gpucb2026} 演变为一套由三阶段构成的融合系统：降维筛选 (Phase 1) $\rightarrow$ 高斯特征代理重建 (Phase 2) $\rightarrow$ 采集-交互闭环 (Phase 3)。



\section{系统框架与描述符特征拆分 (System Architecture \& Feature Descriptor Splitting)}

原始输入数据集 ($X \in \mathbb{R}^{N \times 19}$) 映射到 19 项独立物理化学描述符 (Descriptors)，并覆盖了 CO$_2$ 还原体系的主要特征和结构指标。若强行使用 19 维特征展开代理模型拟合，受制于局部低信噪比与可获取数据点数量，预测面可能会落入过度平滑或局部震荡两个极端。因此，系统实现了分离式的提取和重建机制，并使用“前置填补 (Pre-fill-before-choice)”逻辑予以修复。

\subsection{非线性特征拆分及基尼重要度评判}
有别于线性主成分分析 (PCA) 带来的物理解释性丧失，我们依靠随机森林树群对所有维度特征作相对重要度（Feature Importances）估算，并截断以保留 Top-$K$ 描述符：

\begin{equation}
    \text{Imp}(x_j) = \frac{1}{N_{trees}} \sum_{t=1}^{N_{trees}} \sum_{s \in \text{nodes}(t), v(s) = x_j} p(s) \cdot \Delta I(s)
\end{equation}

$\Delta I(s)$ 为基于特征 $x_j$ 拆解引起的杂质 (通常为方差) 下降，$p(s)$ 为节点权重。此阶段确保最不相关或最受噪声干扰的维度在 GP-BO 前即被丢弃，生成 $X_{selected} \in \mathbb{R}^{N \times K}$，并保证了随后的局部采样有效性\cite{turbo2019}。

\subsection{缺失特征重构与“Pre-fill-before-choice”机制}
剥离了弱关联维度后，算法在获取最新优化点 $x^*_{cont} \in \mathbb{R}^K$ 计算最优收率。然而由于真实合成需全套 19 维配方基准（Strict Feature Schema），若静默抛弃则将直接导致复审阶段化学可行性错位。为此架构引入 \textbf{基于历史映射的前突预填充 (Pre-fill-before-choice)}，对那些处于 $X_{non-selected}$ 的特征，采用数据集同源近似插值或默认基准将其恢复投影至原始 19 维设计域，在呈现给领域专家前保证描述符状态闭环，既发挥了降维优势又实现了信息一致保真。



\section{高斯过程代理模型的数学重构 (Surrogate Model Mathematical Reformulation)}

\subsection{具有独立长度尺度的 ARD Matérn 核}
利用观测对象 $\mathcal{D}_N = \{X_{selected}^{(i)}, y^{(i)}\}_{i=1}^N$ 构建代理模型 $f(X) \sim \mathcal{GP}(\mu, k(X, X'))$。

化学体系内诸如功函数(eV)或者比表面积($\text{m}^2/\text{g}$)的标度差异悬殊，采用通用的欧里几得度量会导致量纲坍塌。报告选择采用带 \textbf{自动关联性决定 (Automatic Relevance Determination, ARD)} 的 Matérn 核，其平滑度设置为 $\nu = 2.5$ 兼顾不可导的高频震荡：
\begin{equation}
    k_{\text{Matérn } 2.5}(x, x') = \sigma_f^2 \left( 1 + \sqrt{5} d_{\text{ARD}} + \frac{5}{3} d_{\text{ARD}}^2 \right) \exp\left( -\sqrt{5} d_{\text{ARD}} \right)
\end{equation}
引入加权对角矩阵，对经过特征筛选存留的 $K$ 维变量实现独立长度尺度缩放 $l_k$：
\begin{equation}
    d_{\text{ARD}}(x, x') = \sqrt{ \sum_{k=1}^K \frac{(x_k - x'_k)^2}{l_k^2} }
\end{equation}
此处 $l_k$ 的自适应有效补偿了特征值之间的刚性错位。对于更为复杂的形貌突变响应曲面，引入复合频谱混合核 (Spectral Mixture Kernels)\cite{catbox2025} 的设计作为可选接口已被预留。

\subsection{物理显式的边界感知归一化 (Bounds-Aware Normalization)}
数据映射至 $[0,1]^K$ 是 BO 后端的常态化要求。但传统机器学习常采用 $\text{MinMax}(X_{train})$，这会在模型外推超出历史样本极大极小值时产生灾难性失真：
系统改用\textbf{理论设计域外推下界归一化}，设外部专家给予已知特征的实验理论边界配置字典 \texttt{DESIGN\_SPACE\_BOUNDS} 为 $(l^{(c)}, u^{(c)})$：
\begin{equation}
    x_{norm, k} = \max \left( 0, \;\min \left( 1, \; \frac{x_k - l^{(k)}}{u^{(k)} - l^{(k)}} \right) \right)
\end{equation}
这强制确立了高斯边缘回归能够无缝探针“未开发区域”，完全杜绝由训练记录极限作为归一化分母引起的伪外推偏误。

\subsection{超参数的边际对数似然估计}
为确定核平稳方差 $\sigma_f^2$，观测白噪声方差 $\sigma_n^2$ 以及尺度向量 $\mathbf{l}$，对高斯模型作边际分布积分，最大化 Exact Marginal Log-Likelihood (MLL)：
\begin{equation}
    \log p(\mathbf{y} | X, \bm{\theta}) = -\frac{1}{2} \mathbf{y}^\top (K(X, X; \bm{\theta}) + \sigma_n^2 I)^{-1} \mathbf{y} - \frac{1}{2} \log|K(X, X; \bm{\theta}) + \sigma_n^2 I| - \frac{N}{2} \log 2\pi
\end{equation}
求解过程使用 L-BFGS-B 或基于 Torch 的梯度下降展开。



\section{采集策略与交互优化 (Acquisition \& Interactive Optimization)}

高分辨收率测定常带有由于光源扰动或杂质吸附带来的固有环境噪声。我们结合问题情况，融合了置信区间探索与带躁期望改进两类算法基准。

\subsection{基于 UCB 的时频调度}
依据 Ding et al.\cite{gpucb2026} 之 GP-UCB 理论准则，最大化置信上边界即可在探索不确定区域(Exploration)与收割当前高收率点(Exploitation)产生黄金折中：
\begin{equation}
    \alpha_{\text{UCB}}(x) = \mu_{t-1}(x) + \beta_t^{1/2} \sigma_{t-1}(x)
\end{equation}
本框架设计支持按迭代轮数动态控制退火机制（`theory` 模式，随着采样深化提升 $\beta_t$ 防止困于局部极值），或是保留用户指定恒定探索参数。

\subsection{噪声干扰下的 Expected Improvement (qNEI)}
为了使得收率优化针对环境干扰有更强鲁棒性\cite{cbo2017}，BoTorch 引擎底层的带噪期望提升
函数 (\texttt{qNoisyExpectedImprovement}) 运用蒙特卡洛法整合观测概率提升：
\begin{equation}
    \alpha_{\text{qNEI}}(x) = \mathbb{E}_{y|x \sim \mathcal{N}}\left[ \max(y - y_{best}^*,\, 0) \right]
\end{equation}
使用 128 采样数目的近似 SobolQMCNormalSampler 对高阶期望作 Quasi-Monte Carlo 无偏估计\cite{botorch2019}。



\section{“人类在环”机制及伪代码实现 (Human-in-the-Loop Implementation)}

在连续边界寻优得到的连续理想配方点 $x^*_{cont}$ 与读取外部历史离散库得到的候选最佳点 $X^*_{disc}$ 之间，模型未必保证它们对应的物化配置绝对可以在实验室合规组装。\textbf{此时人工干预(Manual Override)作为控制流汇聚核心介入}。

\begin{algorithm}[H]
\caption{CO$_2$RR Human-in-the-Loop 贝叶斯优化执行流}
\label{alg:hitl-bo}
\begin{algorithmic}[1]
\REQUIRE 特征边界 $\mathcal{B}$，总迭代数 $T$，初始观测组 $\mathcal{D}_0$，候选数据库 $D_{cand}$
\STATE \textbf{Initialize:} 校验 19 维特征。用户选定产物收率维度 (如 \texttt{Y\_CO}) 以及采集策略。
\STATE \textbf{Phase 1 (Feature Selection)} \textit{// 启发式信息压缩}
\STATE 使用 RF 回归器从 $\mathcal{D}_0$ 对目标组分配权重，排序筛选最显著的 $K$ 项组作 $X_{selected}$。
\FOR{$t = 1 \textbf{ to } T$}
    \STATE 定义设计边界感知归一项变换与标准化函数。
    \STATE $\mathcal{M}_{GP} \leftarrow \text{Fit}(\mathcal{D}_{t-1}, k_{\text{ARD}})$. （见 Section 3）
    \STATE \textbf{Acquisition Maximization} \textit{// L-BFGS-B 重启采样}
    \STATE $x_{cont} \leftarrow \arg\max_{x \in \mathcal{B}_{K}} \alpha_{\text{utility}}(x; \mathcal{M}_{GP})$
    \STATE 从数据库计算剩余组合候选预测 $\hat{Y}_{disc} \leftarrow \text{Evaluate}(D_{cand}, \mathcal{M}_{GP})$，推举 Top-N。
    \STATE \textbf{Human Review \& Pre-fill} \textit{// 人机校验与信息还原}
    \STATE 映射 $x_{cont}$ 返回全局 19 维设计域并输出预测。专家基于化学直觉判断。
    \IF{专家不满意所有由算法提议的连/离散候选项}
        \STATE 执行 \texttt{Manual Override} : 强行注入自行判读的高可行性新配方 $x_{manual}$。
        \STATE $x_{final\_exp} \leftarrow x_{manual}$
    \ELSE
        \STATE $x_{final\_exp} \leftarrow$ 选定的由算法推论的优势项。
    \ENDIF
    \STATE \textbf{Laboratory Execution} : 制备对应催化结构，运行 CO$_2$ 全析出反应，取得多种类收率 $\mathbf{Y}_{all\_products}$。
    \STATE \textbf{Update} : $\mathcal{D}_t \leftarrow \mathcal{D}_{t-1} \cup (x_{final\_exp}, \mathbf{Y}_{all\_products})$。
\ENDFOR
\end{algorithmic}
\end{algorithm}



\section{结语}
结合多维度特征解耦裁剪、带 ARD 核物理感知映射的高斯代理，到基于 GP-UCB 和人类干预逻辑重写的新一轮闭环 BO 迭代模型，构架不仅在数学严格性上保证了 CO$_2$RR 等不确定大搜索面下的收敛界限，也是将当前顶层 ML 理论与底层化学专家知识作最佳弥合的技术典范。该基建将在多成分催化体系和高通量实验室合成间成为有力的驱动支点。

\newpage
\begin{thebibliography}{99}

\bibitem{gpucb2026}
Ding et al., \textit{Efficient and Principled Scientific Discovery through Bayesian Optimization: A Tutorial}. arXiv:2604.01328v3.

\bibitem{cbo2017}
Letham et al., \textit{Constrained Bayesian Optimization with Noisy Experiments}. arXiv:1706.07094v2.

\bibitem{botorch2019}
Balandat et al., \textit{BoTorch: A Framework for Efficient Monte-Carlo Bayesian Optimization}. arXiv:1910.06403v3.

\bibitem{turbo2019}
Wang et al., \textit{Scalable Global Optimization via Local Bayesian Optimization}. arXiv:1910.01739v4.

\bibitem{bo2012}
Snoek et al., \textit{Practical Bayesian Optimization of Machine Learning Algorithms}. arXiv:1206.2944v2.

\bibitem{catbox2025}
Li et al., \textit{CatBOX: A Categorical-Continuous Bayesian Optimization with Spectral Mixture Kernels for Accelerated Catalysis Experiments}. arXiv:2505.17393v2.

\bibitem{jacs2024}
Author et al., \textit{Multivariate Bayesian Optimization of CoO Nanoparticles for CO2 Hydrogenation Catalysis}. Journal of the American Chemical Society (2024), DOI: 10.1021/jacs.4c03789.

\bibitem{nature2025}
Author et al., \textit{Identifying a highly efficient molecular photocatalytic CO2 reduction system via descriptor-based high-throughput screening}. Nature Catalysis (2025), DOI: 10.1038/s41929-025-01291-z.

\end{thebibliography}

\end{document}